This version's system requirements were much more flexible than those imposed in \autoref{sec:v0-requirements}.

\begin{table}[H]
    \centering
    \begin{tabular}{|l|p{10cm}|}
    \hline
    \textbf{No.} & \textbf{Description} \\
    \hline
    1 & Requires standard interfaces to communicate with a flight computer for potential Flatsat integration. \\
    \hline
    2 & Avoid BGA package components whenever possible to reduce production costs. \\
    \hline
    3 & The processor and image sensor must be on independent boards. \\
    \hline
    4 & The image sensor board must use an M12 lens mount. \\
    \hline
    5 & Be able to save image data to a microSD card (unidirectional). \\
    \hline
    6 & Must interface with an CIS via DCMI. \\
    \hline
    \end{tabular}
    \caption{C2Sv0.1 system requirements.}
    \label{tab:v0.1-requirements}
\end{table}

Comparing \autoref{tab:v0.1-requirements} and \autoref{tab:v0-requirements}, we can see that Verion 0 focused on strict technical requirements like pixel count, whereas Version 0.1 uses a high-level systems approach to focus on the must-haves. These relaxed requirements helped focus on developing functional circuits rather than optimising them before I even knew if they worked.
\nl
In addition to changing system requirements, all of the mission requirements were removed from this version as the previous version required as aspects like specific pixel resolution are unnecessary for a prototype.